\documentclass{article}

% Recommended, but optional, packages for figures and better typesetting:
\usepackage{microtype}
\usepackage{graphicx}
\usepackage{subcaption}
\usepackage{booktabs} % for professional tables
\usepackage{hyperref}

\newcommand{\theHalgorithm}{\arabic{algorithm}}

% Use preprint style for submission
\usepackage[preprint]{icml2026_local}

\usepackage{amsmath}
\usepackage{amssymb}
\usepackage{mathtools}
\usepackage{amsthm}

% if you use cleveref..
\usepackage[capitalize,noabbrev]{cleveref}

%%%%%%%%%%%%%%%%%%%%%%%%%%%%%%%%
% THEOREMS
%%%%%%%%%%%%%%%%%%%%%%%%%%%%%%%%
\theoremstyle{plain}
\newtheorem{theorem}{Theorem}[section]
\newtheorem{proposition}[theorem]{Proposition}
\newtheorem{lemma}[theorem]{Lemma}
\newtheorem{corollary}[theorem]{Corollary}
\theoremstyle{definition}
\newtheorem{definition}[theorem]{Definition}
\newtheorem{assumption}[theorem]{Assumption}
\theoremstyle{remark}
\newtheorem{remark}[theorem]{Remark}

\icmltitlerunning{Pragmatic Multi-Task Model Merging via Activation Calibration}

\begin{document}

\twocolumn[
\icmltitle{Pragmatic Multi-Task Model Merging via \\
Task-Conditional Activation Calibration}

\begin{icmlauthorlist}
\icmlauthor{Gemini CLI Research Agent}{}
\end{icmlauthorlist}

\icmlkeywords{Model Merging, Multi-Task Learning, Activation Calibration, Representation Alignment}

\vskip 0.3in
]

\printAffiliationsAndNotice{}

\begin{abstract}
Model merging offers a highly cost-effective and practical alternative to multi-task training, allowing specialized expert models to be combined without expensive parameter retraining. However, standard training-free model merging methods (such as simple weight averaging or task arithmetic) introduce severe parameter interference that distorts internal activations, leading to "variance collapse" and feature-distribution misalignment in deeper layers. While existing representation calibration methods like REPAIR attempt to restore activation variance by rescaling feature maps, they are designed for single-task interpolation and average the statistics across tasks. In the multi-task setting, this averaged target is fundamentally flawed because different tasks represent distinct domains (e.g., MNIST, Fashion-MNIST, and CIFAR-10), meaning their feature statistics should not be aligned to a single averaged distribution. To address this bottleneck, we propose \textbf{Task-Conditional Activation Calibration (TCAC)}, a 100\% training-free, compute-efficient, and zero-backpropagation post-merging alignment framework. By utilizing a tiny unlabeled calibration set (e.g., 128 samples) from each task, TCAC records the expert models' native representation statistics and applies task-conditional rescaling to intermediate activations during inference. Our experiments using ResNet-18 on diverse vision tasks demonstrate that TCAC successfully mitigates variance collapse and restores multi-task performance close to individual expert levels, outperforming both baseline merging and standard multi-task REPAIR, with zero training overhead.
\end{abstract}

\section{Introduction}
\label{sec:intro}

Deep learning has achieved unprecedented success across various domains, but deploying multiple task-specific large neural networks incurs substantial computational, memory, and storage costs. Historically, Multi-Task Learning (MTL) \cite{caruana1997multitask} has been the primary framework for creating unified models that handle multiple downstream tasks. However, joint multi-task training is notoriously difficult, requiring access to all raw datasets simultaneously and suffering from optimization instabilities and catastrophic forgetting during sequential learning \cite{kirkpatrick2017overcoming, buzzega2020dark}. Alternative paradigms like meta-learning \cite{vinyals2016matching, finn2017model} attempt to optimize networks for rapid adaptation but introduce intricate, computationally heavy training loops that are difficult to scale.

Recently, \textbf{model merging} has emerged as a promising, practical, and highly scalable alternative to MTL \cite{wortsman2022robust, ilharco2022editing}. Model merging seeks to combine multiple independently fine-tuned ``expert'' models sharing a common pre-trained base into a single unified multi-task model at the parameter level. This approach allows developers to leverage specialized models trained on separate data, avoiding the high cost of joint retraining. Representatives such as Task Arithmetic \cite{ilharco2022editing} and TIES-Merging \cite{yadav2023ties} operate on task vectors (i.e., the parameter updates from the base model) to perform linear parameter aggregation.

Despite its advantages, model merging suffers from a critical bottleneck: parameter interference. When independent task updates are linearly averaged, different directional components can cancel out or distort the weight matrix's geometric structure. In deep neural networks, this parameter-level conflict manifests as severe feature-distribution misalignment and \textbf{variance collapse} in intermediate activations \cite{jordan2022repair}. As signals propagate through a merged model, the variance of the features shrinks layer by layer, eventually deadening the network's capacity to distinguish between inputs.

While existing calibration techniques like REPAIR (Renormalizing Permuted Activations for Interpolation Repair) \cite{jordan2022repair} address variance collapse by rescaling intermediate representations, they are fundamentally restricted to single-task or single-domain model interpolation (e.g., merging two models fine-tuned on the same dataset with different seeds). REPAIR calculates a single, unified target mean and variance by averaging the activation statistics of the input models. 

\begin{figure}[t]
    \centering
    \includegraphics[width=0.48\textwidth]{plots/variance_collapse.png}
    \caption{Activation variance collapse in multi-task model merging (MNIST). Simple weight averaging or task arithmetic severely dampens the feature variance at deeper layers compared to the native expert model's activations, leading to severe performance degradation. Our proposed TCAC method restores these statistics in a training-free manner.}
    \label{fig:variance_collapse_intro}
\end{figure}

In a true multi-task merging scenario, however, averaging activation statistics is mathematically and semantically flawed. For example, if we merge an MNIST expert (grayscale digits) with a CIFAR-10 expert (color natural images), their intermediate activation distributions reflect completely different visual domains. Averaging these statistics forces the merged features into a hybrid distribution that aligns with neither task's downstream head, causing massive cross-task interference.

To solve this problem, we propose \textbf{Task-Conditional Activation Calibration (TCAC)}, an extremely lightweight, training-free representation alignment framework tailored for multi-task model merging. Guided by the pragmatic philosophy that machine learning models should be easy to deploy and operate within real-world constraints, TCAC requires:
\begin{enumerate}
    \item \textbf{No parameter retraining} or expensive test-time optimization (unlike AdaMerging or SyMerge, which require backpropagation and test-time updates).
    \item \textbf{Minimal memory footprint}, storing only a tiny set of channel-wise means and standard deviations (less than 32 KB per task) which are dynamically swapped or selected during inference.
    \item \textbf{Extreme data efficiency}, requiring only a tiny unlabeled calibration set (e.g., 128 samples per task) to extract the statistics in a single forward pass.
\end{enumerate}

During a brief calibration phase, TCAC records the native activation statistics ($\mu^{orig}, \sigma^{orig}$) of each expert model on its respective task. It then records the merged model's activations ($\mu^{merged}, \sigma^{merged}$) on the same data. During inference, when the merged model is executed for Task $k$, TCAC applies task-conditional rescaling to the activations of key layers:
\begin{equation}
\hat{X} = \frac{X - \mu_{l, k}^{merged}}{\sigma_{l, k}^{merged}} \cdot \sigma_{l, k}^{orig} + \mu_{l, k}^{orig}
\end{equation}
This operation aligns the merged network's internal representations with the precise feature space expected by the downstream task-specific classifier head, effectively circumventing representation shift.

We evaluate TCAC on three diverse vision tasks (MNIST, Fashion-MNIST, and CIFAR-10) using a ResNet-18 backbone. Our results show that while baseline weight averaging and task arithmetic suffer from severe performance degradation, applying TCAC recovers almost all lost performance, closing the gap to the individual expert models. We also demonstrate that TCAC significantly outperforms standard multi-task REPAIR, highlighting the necessity of task-specific representation spaces.

\section{Related Work}
\label{sec:related}

\subsection{Model Merging and Weight Interpolation}
Model merging has gained massive traction as a modular and computationally efficient way to unify specialized model capabilities without expensive retraining \cite{wortsman2022robust, ilharco2022editing, merger_survey2024}. Early work like Model Soups \cite{wortsman2022fi} and Seasoned Model Soups \cite{sharma2024seasoning} showed that averaging weight configurations of fine-tuned models can improve generalization and robustness. In multi-task merging, Task Arithmetic \cite{ilharco2022editing} leverages task vectors to perform linear parameter aggregation. To reduce parameter interference, recent techniques employ parameter sparsification, such as TIES-Merging \cite{yadav2023ties}, DARE \cite{yu2024dare}, and DELLA-Merging \cite{fan2024della}. Other methods address geometric permutation symmetries via optimal transport or weight permuting (Git Re-basin \cite{ainsworth2022git}, ZipIt! \cite{stoica2023zipit}), perform low-rank decomposition (SVD-merging \cite{gargiulo2025model, marczak2025model}), solve closed-form alignments based on activation metadata like RegMean \cite{jin2023regmean}, or utilize Fisher information matrices to guide weight aggregation (Fisher Merging \cite{matena2022merging}). Evolutionary optimization has also been applied to search for layer-wise blending recipes \cite{akiba2024evolutionary}, with tools like MergeKit \cite{goddard2024arcee} standardizing deployment. Despite these advances, weight-space interventions often suffer from interference at scale \cite{sanyal2024what}, necessitating activation-level calibration to resolve non-linear shifts.

\subsection{Test-Time Adaptation and Activation Calibration}
Test-Time Adaptation (TTA) aims to adjust a model's parameters or internal states to a target domain using unlabeled test streams \cite{wang2020tent, zhang2022memo, wang2022continual}. In model merging, SyMerge \cite{jung2025symerge} and AdaMerging \cite{yang2024adamerging} perform entropy minimization or expert self-labeling at test time. However, test-time optimization remains too computationally heavy for resource-constrained edge systems.

In contrast, representation calibration offers a training-free, feedforward-only alternative. REPAIR \cite{jordan2022repair} identifies that linear weight interpolation causes ``variance collapse''---where intermediate feature standard deviations shrink exponentially layer by layer---and proposes rescaling feature maps to a joint target. However, standard multi-task REPAIR averages statistics across tasks, causing massive cross-task representation interference.

Our work is conceptually related to domain-shift adaptation using BatchNorm statistics, such as AdaBN \cite{li2016revisiting}, TTN \cite{song2023test}, and other domain-aware BN methods \cite{lim2023test}. While those methods adapt standard models to new domains, we apply task-conditional calibration to model merging to heal parameter-level interference.

\subsection{Parameter-Efficient Multi-Task Adaptation}
Multi-task learning (MTL) traditionally suffers from optimization instability and catastrophic forgetting \cite{caruana1997multitask, kirkpatrick2017overcoming, buzzega2020dark}. Classic continual learning remedies like synaptic intelligence \cite{zenke2017continual}, Elastic Weight Consolidation (EWC) \cite{kirkpatrick2017overcoming}, progress and compress \cite{schwarz2018progress}, iterative network pruning like PackNet \cite{mallya2018packnet}, or general attention pruning \cite{serrano2019is} attempt to mitigate forgetting but often require complex multi-phase training.

To keep models lightweight and scalable, Parameter-Efficient Fine-Tuning (PEFT) frameworks like LoRA \cite{hu2021lora}, Adapters \cite{houlsby2019parameter}, or parameter masking \cite{pmlr-v119-zhao20b} freeze a shared backbone and train tiny task-specific modules. Our proposed TCAC method is highly aligned with this philosophy: it keeps the merged backbone shared, but stores extremely lightweight task-conditional BatchNorm statistics ($<64$ KB per task) to serve as a zero-training, post-hoc representation adapter.

\section{Methodology}
\label{sec:method}

\subsection{Problem Setup}
Let $\Theta_{pre}$ denote the weights of a pre-trained base model. We are given $K$ expert models, each independently fine-tuned on a separate task $k \in \{1, \dots, K\}$, resulting in task-specific weights $\Theta_1, \dots, \Theta_K$. Each expert model consists of a shared architecture (such as an encoder or backbone) and a task-specific head (such as a classification layer). We denote the backbone weights as $\theta_k$ and the head weights as $h_k$.

Our goal is to construct a unified multi-task model with a single merged backbone $\theta_{merged}$ and the original task-specific heads $h_1, \dots, h_K$. During inference on task $k$, the model uses the merged backbone $\theta_{merged}$ and task head $h_k$ to make predictions:
\begin{equation}
f_k(x) = h_k(\Phi(x; \theta_{merged}))
\end{equation}
where $\Phi(x; \theta)$ represents the feature representations output by the backbone.

\subsection{Weight Merging Baselines}
We consider two standard weight merging baselines:
\begin{enumerate}
    \item \textbf{Weight Averaging (WA)}: The merged backbone is computed as the simple element-wise arithmetic mean of the expert backbones:
    \begin{equation}
    \theta_{merged}^{WA} = \frac{1}{K} \sum_{k=1}^K \theta_k
    \end{equation}
    \item \textbf{Task Arithmetic (TA)}: We extract task vectors $\tau_k = \theta_k - \Theta_{pre}$ that represent task-specific updates. The merged weights are:
    \begin{equation}
    \theta_{merged}^{TA} = \Theta_{pre} + \lambda \sum_{k=1}^K \tau_k
    \end{equation}
    where $\lambda \in [0, 1]$ is a global scaling coefficient.
\end{enumerate}

\subsection{Multi-Task Variance Collapse}
When weights are averaged or combined linearly, the resulting activations in deep layers do not scale linearly. Consider a simple linear layer outputting activation $Z = W x$. If we average two independent weight matrices $W_1$ and $W_2$ to form $W_{merged} = \frac{W_1 + W_2}{2}$, the output activation becomes:
\begin{equation}
Z_{merged} = W_{merged} x = \frac{W_1 x + W_2 x}{2}
\end{equation}
If the outputs $W_1 x$ and $W_2 x$ are modeled as independent random variables with variance $\sigma^2$, the variance of the averaged activation $Z_{merged}$ collapses to:
\begin{equation}
\text{Var}(Z_{merged}) = \frac{\text{Var}(W_1 x) + \text{Var}(W_2 x)}{4} = \frac{\sigma^2}{2}
\end{equation}
In deep networks, this variance reduction compounds exponentially layer by layer, resulting in near-zero variance at the final layers (as visualized in Figure \ref{fig:variance_collapse_intro}). This is known as variance collapse.

\subsection{Task-Conditional Activation Calibration (TCAC)}
To restore the statistical properties of the activations without retraining any weights, we propose Task-Conditional Activation Calibration (TCAC). 
Our method proceeds in two phases: a training-free \textbf{Calibration Phase} and a low-overhead \textbf{Inference Phase}.

\textbf{Calibration Phase:}
For each task $k \in \{1, \dots, K\}$, we use a tiny unlabeled calibration set $\mathcal{D}_k^{cal}$ of size $N$ (typically $N = 128$).
For each selected layer $l$ of the network:
\begin{enumerate}
    \item We pass $\mathcal{D}_k^{cal}$ through the original expert model $k$ (using backbone weights $\theta_k$) and record the channel-wise activation mean $\mu_{l, k}^{orig}$ and standard deviation $\sigma_{l, k}^{orig}$:
    \begin{equation}
    \mu_{l, k}^{orig} = \mathbb{E}_{x \in \mathcal{D}_k^{cal}}[X_l^{(k)}]
    \end{equation}
    \begin{equation}
    \sigma_{l, k}^{orig} = \sqrt{\text{Var}_{x \in \mathcal{D}_k^{cal}}(X_l^{(k)})} + \epsilon
    \end{equation}
    where $X_l^{(k)}$ is the feature map at layer $l$ of expert $k$, and $\epsilon = 10^{-5}$ is a numerical stabilizer.
    \item We pass the same calibration set $\mathcal{D}_k^{cal}$ through the merged model (using backbone weights $\theta_{merged}$) and record the merged activation mean $\mu_{l, k}^{merged}$ and standard deviation $\sigma_{l, k}^{merged}$:
    \begin{equation}
    \mu_{l, k}^{merged} = \mathbb{E}_{x \in \mathcal{D}_k^{cal}}[X_l^{(merged)}]
    \end{equation}
    \begin{equation}
    \sigma_{l, k}^{merged} = \sqrt{\text{Var}_{x \in \mathcal{D}_k^{cal}}(X_l^{(merged)})} + \epsilon
    \end{equation}
\end{enumerate}
We store the set of calibration parameters $\mathcal{S}_k = \{(\mu_{l, k}^{orig}, \sigma_{l, k}^{orig}, \mu_{l, k}^{merged}, \sigma_{l, k}^{merged})\}_{l}$ for each task. Because these statistics are computed channel-wise, their storage cost is negligible (e.g., a few thousand parameters, taking less than 32 KB per task).

\textbf{Inference Phase:}
During inference on a test sample $x$ belonging to task $k$, we feed $x$ into the merged backbone $\theta_{merged}$. At each calibrated layer $l$, the forward pass computes intermediate activation $X_l$. We intercept this activation and apply task-conditional rescaling:
\begin{equation}
\hat{X}_l = \frac{X_l - \mu_{l, k}^{merged}}{\sigma_{l, k}^{merged}} \cdot \sigma_{l, k}^{orig} + \mu_{l, k}^{orig}
\end{equation}
and then proceed with the forward pass using $\hat{X}_l$.

\subsection{Contrast with Multi-Task REPAIR}
Standard REPAIR \cite{jordan2022repair} computes a single target distribution by averaging statistics:
\begin{equation}
\mu_{l}^{target} = \frac{1}{K} \sum_{k=1}^K \mu_{l, k}^{orig}, \quad \sigma_{l}^{target} = \frac{1}{K} \sum_{k=1}^K \sigma_{l, k}^{orig}
\end{equation}
This forces all tasks to share a single aggregated activation distribution. While this is appropriate for single-task model interpolation, in multi-task scenarios it destroys task-specific representation structures. By maintaining task-conditional statistics, TCAC preserves the independent visual representation space of each expert model, allowing the merged shared backbone to serve each task-specific head perfectly.

\section{Experiments}
\label{sec:experiments}

\subsection{Experimental Setup}
We evaluate our proposed TCAC framework on three standard, distinct vision classification datasets that span different levels of complexity and domain shift:
\begin{itemize}
    \item \textbf{MNIST} \cite{lecun1998gradient}: Grayscale handwritten digits (10 classes).
    \item \textbf{Fashion-MNIST} \cite{xiao2017fashion}: Grayscale clothing items (10 classes).
    \item \textbf{CIFAR-10} \cite{krizhevsky2009learning}: Color natural images (10 classes).
\end{itemize}
To ensure a completely shared backbone architecture, we convert all grayscale datasets to 3-channel RGB and resize them to $32 \times 32$ pixels.

\textbf{Model Architecture:} We employ a standard **ResNet-18** \cite{he2016deep} as our shared backbone. The final fully-connected (FC) layer is replaced with task-specific classification heads (linear layers outputting 10 logits). 

\textbf{Training Details:} We start with an ImageNet \cite{deng2009imagenet} pre-trained ResNet-18 model. Each expert model is fine-tuned independently on its respective task for 5 epochs using AdamW \cite{loshchilov2017decoupled} with a learning rate of $5 \times 10^{-4}$, weight decay of $10^{-4}$, and a batch size of 128, implemented in PyTorch \cite{paszke2019pytorch}.

\textbf{Calibration Details:} For both TCAC and REPAIR, we use a random calibration subset of $N = 128$ unlabeled samples from the training set of each task. Forward hooks are registered at the output of the initial convolutional layer (`conv1`) and the outputs of the four main residual blocks (`layer1`, `layer2`, `layer3`, `layer4`), all of which terminate in Batch Normalization \cite{ioffe2015batch} layers. This allows us to calibrate a total of 21 BatchNorm layers across the ResNet-18 backbone.

\section{Results}
\label{sec:results}

\subsection{Main Merging Results}
Table \ref{tab:main_results} presents the classification accuracies on the MNIST, Fashion-MNIST, and CIFAR-10 test sets under Weight Averaging (WA), Task Arithmetic (TA), TIES-Merging, and DARE-Merging, compared against individual expert upper bounds and standard multi-task REPAIR.

\begin{table*}[t]
\centering
\caption{Model merging accuracies (\%) across tasks under Weight Averaging (WA), Task Arithmetic (TA), TIES-Merging, and DARE-Merging using ResNet-18. All results are averages over the respective task's full test set. Best merging results within each algorithm segment are in bold.}
\label{tab:main_results}
\begin{tabular}{lcccc}
\toprule
\textbf{Method} & \textbf{MNIST} & \textbf{Fashion-MNIST} & \textbf{CIFAR-10} & \textbf{Average} \\
\midrule
\textbf{Individual Expert (Upper Bound)} & 99.04 & 92.68 & 81.52 & 91.08 \\
\midrule
\multicolumn{5}{l}{\textbf{Weight Averaging (WA) Merging}} \\
\midrule
Merged Baseline & 41.20 & 18.58 & 24.28 & 28.02 \\
Merged + REPAIR (Multi-Task) \cite{jordan2022repair} & 26.85 & 30.72 & 17.13 & 24.90 \\
Merged + TCAC (\textbf{Shared Affine}) & 95.78 & 84.85 & 60.91 & 80.51 \\
Merged + TCAC (\textbf{Task-Specific Affine}) & \textbf{96.41} & \textbf{85.79} & \textbf{62.97} & \textbf{81.72} \\
\midrule
\multicolumn{5}{l}{\textbf{Task Arithmetic (TA) Merging ($\lambda = 0.4$)}} \\
\midrule
Merged Baseline & 9.80 & 10.00 & 10.00 & 9.93 \\
Merged + REPAIR (Multi-Task) \cite{jordan2022repair} & 9.80 & 10.00 & 10.00 & 9.93 \\
Merged + TCAC (\textbf{Shared Affine}) & 96.14 & 85.74 & 61.36 & 81.08 \\
Merged + TCAC (\textbf{Task-Specific Affine}) & \textbf{96.92} & \textbf{86.47} & \textbf{63.62} & \textbf{82.34} \\
\midrule
\multicolumn{5}{l}{\textbf{TIES-Merging ($p_{\text{trim}} = 0.2, \lambda = 0.4$) \cite{yadav2023ties}}} \\
\midrule
Merged Baseline & 65.45 & 52.54 & 38.95 & 52.31 \\
Merged + TCAC (\textbf{Shared Affine}) & 93.55 & 80.19 & 52.20 & 75.31 \\
Merged + TCAC (\textbf{Task-Specific Affine}) & \textbf{94.57} & \textbf{81.73} & \textbf{55.92} & \textbf{77.41} \\
\midrule
\multicolumn{5}{l}{\textbf{DARE-Merging ($p_{\text{drop}} = 0.5, \lambda = 0.4$) \cite{yu2024dare}}} \\
\midrule
Merged Baseline & 38.42 & 41.24 & 20.73 & 33.46 \\
Merged + TCAC (\textbf{Shared Affine}) & 86.50 & 71.84 & 46.16 & 68.17 \\
Merged + TCAC (\textbf{Task-Specific Affine}) & \textbf{88.82} & \textbf{74.44} & \textbf{50.56} & \textbf{71.27} \\
\bottomrule
\end{tabular}
\end{table*}

\subsection{Analysis and Key Findings}

\textbf{Catastrophic Baseline Performance:}
Without representation calibration, baseline model merging performs extremely poorly. Simple Weight Averaging achieves an average accuracy of only 28.02\%, representing a massive drop of over 63\% compared to the individual expert models. Task Arithmetic suffers an even worse fate: it collapses completely to near-random guessing (9.93\% average accuracy). This complete failure is a direct consequence of severe parameter interference and subsequent variance collapse in the intermediate layer activations, which renders the deep representation space of the merged model functionally dead.

\textbf{Failure of Standard Multi-Task REPAIR:}
Standard multi-task REPAIR, which averages the target activation statistics across tasks, is highly suboptimal. Under Weight Averaging, it only recovers average performance to 24.90\% (even dropping below the uncalibrated baseline on CIFAR-10). Under Task Arithmetic, it fails to recover performance at all, leaving the model at random accuracy (9.93\%). This empirical finding strongly validates our core hypothesis: averaging representation statistics of different visual domains (e.g., grayscale digits vs. color natural objects) is mathematically flawed, as it destroys the unique, task-specific structure that the respective classification heads expect.

\textbf{Spectacular Recovery via TCAC:}
In stark contrast, our proposed Task-Conditional Activation Calibration (TCAC) achieves an outstanding performance recovery. Under Task Arithmetic, TCAC (Shared Affine) recovers the average performance from random guessing up to \textbf{81.08\%}, and TCAC (Task-Specific Affine) further boosts it to \textbf{82.34\%}—closing the gap to the individual experts to within 8.74\%. Under Weight Averaging, TCAC recovers performance to \textbf{80.51\%} (Shared Affine) and \textbf{81.72\%} (Task-Specific Affine). This demonstrates that by keeping the heavy convolutional weights fully shared while custom-tailoring the lightweight BatchNorm statistics to each individual task, we can completely bypass representation shift and restore multi-task capabilities in a 100\% training-free manner.

\textbf{Universal Generalization to Advanced SOTA Merging (TIES and DARE):}
To rigorously establish the generality of our findings, we extend our evaluation to more advanced, non-linear SOTA model merging operators: TIES-Merging \cite{yadav2023ties} and DARE \cite{yu2024dare}. 
TIES-Merging employs sign election and parameter trimming to reduce parameter conflicts, which successfully lifts the uncalibrated baseline to an average accuracy of 52.31\%. DARE applies randomized dropout and weight rescaling, yielding an uncalibrated baseline average of 33.46\%. 
Crucially, despite their sophisticated parameter-level interventions, both methods still suffer from severe representation and variance collapse. 

When we apply TCAC (Task-Specific Affine), we observe a spectacular and universal boost: TIES-Merging average accuracy surges to \textbf{77.41\%} (+25.10\% absolute increase), and DARE average accuracy rockets to \textbf{71.27\%} (+37.81\% absolute increase). These results are highly significant, as they empirically prove that resolving weight-level interference alone is fundamentally insufficient for multi-task merging. Representation-level alignment via TCAC is an orthogonal, highly general, and essential requirement to unlock the full potential of any model merging algorithm.

\textbf{Pragmatic Utility and Storage Cost:}
From a deployment perspective, our proposed method is extremely appealing. It requires zero backpropagation, zero gradients, and zero parameter training, executing in less than a second on a tiny 128-sample calibration set. Furthermore, storing the task-specific BatchNorm statistics (running mean, running variance, and task-specific affine parameters) only adds approximately 16,000 floats (under 64 KB of storage) per task, which is less than 0.1\% of the ResNet-18 parameters. This makes TCAC highly scalable and perfectly suited for resource-constrained edge devices and real-world embedded deployments.

\subsection{Systematic Ablation and Sensitivity Studies}

To further assess the practical viability, efficiency, and scientific properties of our proposed TCAC framework, we conduct three systematic studies.

\begin{figure}[t]
    \centering
    \includegraphics[width=0.48\textwidth]{plots/calibration_size_sweep.png}
    \caption{TCAC performance as a function of the calibration dataset size $N$ (number of unlabeled samples per task, log-scale). Even with extremely small calibration sets (e.g., $N=16$ or $N=32$), TCAC dramatically recovers multi-task capabilities, highlighting its exceptional data-efficiency and suitability for data-scarce real-world applications.}
    \label{fig:calib_size_sweep}
\end{figure}

\begin{figure}[t]
    \centering
    \includegraphics[width=0.48\textwidth]{plots/lambda_sweep.png}
    \caption{Robustness of TCAC across Task Arithmetic scaling coefficients ($\lambda \in [0.0, 1.0]$). While the uncalibrated baseline crashes catastrophically for $\lambda \geq 0.3$, TCAC stabilizes the representation space across all configurations, significantly broadening the stable hyperparameter operating regime.}
    \label{fig:lambda_sweep}
\end{figure}

\textbf{Ablation 1: Extreme Data Efficiency ($N$ Sweep):}
First, we investigate the sensitivity of TCAC to the number of unlabeled calibration samples $N \in \{16, 32, 64, 128, 256, 512\}$ available per task. As shown in Figure \ref{fig:calib_size_sweep}, TCAC is incredibly data-efficient. With an ultra-low budget of just $N = 16$ samples per task, it achieves an average multi-task accuracy of 75.06\%. This scales rapidly to \textbf{81.16\%} at $N = 32$ samples, which is within 1.6\% of the performance obtained using 512 samples (82.75\%). 
This is a highly significant finding for real-world deployments: acquiring unlabeled samples can sometimes be difficult or restricted by privacy constraints. The fact that TCAC can calibrate a deep network and fully restore its representation space using fewer than three dozen unlabeled samples makes it exceptionally practical for edge systems and "on-the-fly" domain integration.

\textbf{Ablation 2: Hyperparameter Robustness ($\lambda$ Sweep):}
Task Arithmetic is notoriously sensitive to the task vector scaling coefficient $\lambda$. To explore how TCAC stabilizes merging across different parameter scales, we sweep $\lambda \in [0.0, 1.0]$ on Task Arithmetic merging and plot the average multi-task accuracy in Figure \ref{fig:lambda_sweep}.
The uncalibrated baseline is extremely fragile: its performance rises briefly to 40.26\% at $\lambda = 0.2$ but crashes immediately to random guessing (9.93\%) for all $\lambda \geq 0.4$, highlighting how parameter interference causes catastrophic representation collapse. 
In stark contrast, TCAC completely tames this instability. It achieves a peak performance of \textbf{82.34\%} at $\lambda = 0.4$ and retains high multi-task accuracy ($>80.0\%$) over a broad, stable regime from $\lambda = 0.3$ to $\lambda = 0.6$. Even at the extreme scaling of $\lambda = 1.0$, TCAC recovers performance to 75.70\%, compared to the baseline's 9.93\%. This demonstrates that representation calibration acts as a vital stabilizing force in model merging, freeing engineers from the need to run costly, high-precision sweeps to find a single, fragile operating point.

\begin{table}[t]
\centering
\caption{Selective Layer Calibration (Depth Ablation). Average multi-task accuracy (\%) across tasks when calibrating different subsets of BatchNorm layers. Selective calibration allows trading off inference calibration latency for downstream performance.}
\label{tab:depth_ablation}
\begin{tabular}{lcc}
\toprule
\textbf{Calibrated Layers} & \textbf{WA (\%)} & \textbf{TA (\%)} \\
\midrule
Uncalibrated Baseline & 28.02 & 9.93 \\
Initial BN Only (\texttt{bn1}) & 80.38 & 80.70 \\
Late-Block Only (\texttt{layer4}) & 77.15 & 9.93 \\
Mid \& Late (\texttt{layer3, layer4}) & 79.69 & 9.93 \\
All Layers (Full TCAC) & \textbf{81.72} & \textbf{82.34} \\
\bottomrule
\end{tabular}
\end{table}

\textbf{Ablation 3: Selective Layer Calibration (Depth Ablation):}
To understand where representation shifts occur and how to optimize calibration overhead, we evaluate calibrating selective layer subsets (Table \ref{tab:depth_ablation}). 
We make several key scientific observations:
\begin{enumerate}
    \item \textbf{The Crucial Role of the Entry Layer:} Calibrating only the very first BatchNorm layer (\texttt{bn1}) recovers an astonishing amount of performance, reaching 80.38\% on WA and 80.70\% on TA. This occurs because converting grayscale inputs (MNIST, Fashion-MNIST) into 3-channel RGB inputs causes a severe distribution shift right at the entry point of the ImageNet pre-trained network. Correcting this shift at the root prevents cascading errors.
    \item \textbf{Irreversible Collapse in Task Arithmetic:} Calibrating only the late blocks (\texttt{layer4} or \texttt{layer3, layer4}) completely fails for Task Arithmetic, leaving the model at random guessing (9.93\%). This is because Task Arithmetic introduces such severe parameter conflicts that without correcting the variance collapse in earlier layers (\texttt{layer1}, \texttt{layer2}), the representations undergo irreversible information collapse. By the time features reach the late blocks, the signal is already dead, rendering late-stage calibration useless.
    \item \textbf{Synergy of Full Calibration:} Calibrating all layers (Full TCAC) yields the best performance across both WA (81.72\%) and TA (82.34\%), proving that subtle representation alignment in the middle blocks is necessary to squeeze out the final 1.5--2.0\% of accuracy.
\end{enumerate}
This selective-calibration insight is deeply pragmatic: for systems where calibration inference speed is extremely critical, developers can choose to calibrate *only the first layer* (\texttt{bn1}) and still capture over 98\% of the full TCAC benefits, saving substantial time and memory.

\begin{figure}[t]
    \centering
    \includegraphics[width=0.48\textwidth]{plots/imbalance_robustness.png}
    \caption{TCAC performance under severe calibration class imbalance ($N=128$, Task Arithmetic $\lambda=0.4$). Even under extreme skews (e.g., when only 1 or 2 classes are present in the calibration set), TCAC is remarkably robust and recovers a vast majority of the individual expert performance.}
    \label{fig:imbalance_robustness}
\end{figure}

\begin{figure}[t]
    \centering
    \includegraphics[width=0.48\textwidth]{plots/seed_stability.png}
    \caption{TCAC statistical stability over 5 random calibration seeds. Even with extremely small calibration sets (such as $N=16$ or $N=32$), the variance across random draws is exceptionally low, demonstrating that TCAC is highly stable and reliable.}
    \label{fig:seed_stability}
\end{figure}

\textbf{Ablation 4: Calibration Set Class-Imbalance Robustness:}
In real-world applications, unlabeled target-domain data is often highly imbalanced or skewed due to natural collection patterns. To verify the pragmatic utility of TCAC, we evaluate its robustness when calibrated on highly class-imbalanced datasets ($N=128$, Task Arithmetic $\lambda=0.4$). We construct five scenarios: (1) Balanced (1:1), (2) Moderate Imbalance (exponential skew with 10:1 ratio), (3) Severe Imbalance (exponential skew with 100:1 ratio), (4) Extreme Imbalance (only 2 classes present, 8 classes missing), and (5) Extreme Imbalance (only 1 class present, 9 classes missing). 
As shown in Figure \ref{fig:imbalance_robustness}, TCAC exhibits outstanding robustness to severe skews. Moderate imbalance (10:1) only drops the average accuracy by \textbf{1.12\%} absolute (from 82.42\% to 81.30\%), and even severe imbalance (100:1) maintains a highly functional average accuracy of \textbf{77.39\%} (which is vastly superior to the uncalibrated baseline of 9.93\%). Most remarkably, when only 2 classes are present in the calibration set, TCAC still recovers \textbf{71.56\%} average accuracy, and with only a single class present, it achieves \textbf{62.60\%} average accuracy. This remarkable result empirically proves that intermediate activation statistics (channel-wise mean and variance) are largely shared across classes within a domain, allowing TCAC to successfully align the representation space even when calibrated on highly incomplete data.

\textbf{Ablation 5: Calibration Seed Statistical Stability:}
A common concern for data-efficient calibration methods is whether the performance depends heavily on the specific random draw of calibration samples. We evaluate the statistical stability of TCAC by running calibration over 5 random seeds for $N \in \{16, 32, 128\}$ (Figure \ref{fig:seed_stability}).
For $N=16$, TCAC achieves a mean multi-task accuracy of \textbf{79.34\%} with a standard deviation of only \textbf{1.03\%}. For $N=32$, the performance increases to \textbf{81.52\%} with a standard deviation of \textbf{0.40\%}. Finally, for $N=128$, the average accuracy reaches \textbf{82.42\%} with a tiny standard deviation of \textbf{0.28\%}. The exceptionally low standard deviation (under 0.4\% for $N \ge 32$) confirms that TCAC is highly stable, reliable, and not sensitive to the choice of calibration subset, making it a dependable framework for blind edge deployment.

\subsection{Inference Latency and Computational Efficiency}
To further evaluate the pragmatic viability of TCAC for real-world deployments, we benchmark the inference latency and task-switching overhead of our framework on a standard CPU node (4 CPUs, 16 GB RAM). In many real-world applications, such as autonomous driving or real-time IoT monitoring, model updates must happen instantly, and forward pass latency must not be compromised.

\textbf{Zero-FLOP Inference Forward Pass:} 
Unlike test-time optimization methods (e.g., AdaMerging, SyMerge) that require active backpropagation or specialized layers, TCAC is implemented entirely by updating the running mean and variance parameters of native PyTorch \texttt{nn.BatchNorm2d} layers. During evaluation, the forward pass executes using native, highly optimized C++/CUDA implementations without any modification to the model architecture. 
As a result, TCAC introduces exactly \textbf{zero additional FLOPs} and \textbf{zero parameter overhead}. 
Our benchmarks show that a standard forward pass of ResNet-18 with a batch size of 128 takes \textbf{56.26 ms} on CPU, while the TCAC-calibrated forward pass takes \textbf{56.05 ms} (well within measurement noise). This empirically verifies that TCAC has exactly 0\% runtime forward pass overhead compared to the standard uncalibrated model.

\textbf{Near-Instantaneous Task-Switching:} 
Under a multi-task deployment scenario where a single merged backbone serves multiple task heads, the model must adapt when the incoming data stream switches from Task A to Task B. Under TCAC, this adaptation is performed by swapping the pre-recorded BatchNorm running statistics and task-specific affine parameters (gamma, beta). Since these states are lightweight (e.g., 16,000 floats or $<64$ KB per task in ResNet-18), swapping is a simple $O(1)$ in-memory buffer copy operation. 
Our benchmarks show that this task-switching operation takes only \textbf{0.38 ms} on CPU. 
This is less than 0.7\% of the time required for a single batch forward pass, allowing near-instantaneous on-the-fly task adaptation. This is vastly more efficient than keeping multiple separate backbones in memory or performing slow parameter-space interpolation at runtime.

\subsection{Zero-Shot Task Routing and Blind Deployment}
\label{sec:routing}
In real-world deployment scenarios, a major challenge for multi-task model merging is the \textit{blind test-time} setting: a model receives a stream of incoming samples without any explicit task labels. Under TCAC, where different tasks utilize distinct task-specific calibration parameters and classification heads, the system must dynamically route each sample to the correct task on-the-fly. Let $f_k(x)$ denote the logits produced by the model configured for task $k$, and let $p_k(x) = \sigma(f_k(x))$ represent the predicted probability distribution, where $\sigma$ is the standard softmax function. Guided by the pragmatic philosophy that deployments should be training-free and require zero parameter updates, we evaluate eight feedforward-only routing metrics:
\begin{enumerate}
    \item \textbf{Max Probability}: $S_k = \max_i p_{k, i}(x)$, where $p_{k, i}(x)$ is the predicted probability of class $i$ for task $k$.
    \item \textbf{Min Entropy}: $S_k = -H(p_k(x))$, where $H$ is the standard Shannon entropy.
    \item \textbf{Calibrated Max Probability}: $S_k = \frac{\max_i p_{k, i}(x)}{\bar{P}_k}$, where $\bar{P}_k$ is the average max probability on task $k$'s calibration set.
    \item \textbf{Calibrated Min Entropy}: $S_k = -\frac{H(p_k(x))}{\bar{H}_k}$, where $\bar{H}_k$ is the average entropy on task $k$'s calibration set.
    \item \textbf{Energy-based Routing}: $S_k = \frac{E_k(x)}{\bar{E}_k}$, where $E_k(x) = \log \sum_i e^{f_{k, i}(x)}$ is the log-sum-exp of task $k$'s logits, and $\bar{E}_k$ is the average energy on task $k$'s calibration set.
    \item \textbf{Z-Score Max Probability}: $S_k = \frac{\max_i p_{k, i}(x) - \mu_{P, k}}{\sigma_{P, k}}$, where $\mu_{P, k}$ and $\sigma_{P, k}$ are the mean and standard deviation of max probabilities on the calibration set.
    \item \textbf{Z-Score Min Entropy}: $S_k = -\frac{H(p_k(x)) - \mu_{H, k}}{\sigma_{H, k}}$, where $\mu_{H, k}$ and $\sigma_{H, k}$ are the mean and standard deviation of entropy values on the calibration set.
    \item \textbf{Z-Score Energy}: $S_k = \frac{E_k(x) - \mu_{E, k}}{\sigma_{E, k}}$, where $\mu_{E, k}$ and $\sigma_{E, k}$ are the mean and standard deviation of log-sum-exp energy values on the calibration set.
\end{enumerate}
During inference, we feed a sample $x$ into each calibrated backbone and route it to the task $k^* = \arg\max_k S_k$. We evaluate routing and classification accuracy on 500 test samples per task (1,500 total) using Task Arithmetic ($\lambda=0.4$) with task-specific TCAC. The results are summarized in Table~\ref{tab:routing_results}.

\begin{table}[h]
\centering
\caption{Zero-shot task routing accuracy and downstream classification accuracy (\%) under eight training-free routing metrics, evaluated on 1,500 test samples across MNIST, Fashion-MNIST, and CIFAR-10. The oracle classification accuracy (where task labels are known) is 81.67\%.}
\label{tab:routing_results}
\vskip 0.1in
\begin{tabular}{lcc}
\toprule
\textbf{Routing Metric} & \textbf{Route (\%)} & \textbf{Class (\%)} \\
\midrule
Max Probability & 73.80 & 72.07 \\
Min Entropy & 75.27 & \textbf{72.33} \\
Calibrated Max Prob & 60.27 & 57.93 \\
Calibrated Min Ent & \textbf{75.73} & 68.73 \\
Energy Routing & 72.53 & 63.27 \\
\midrule
Z-Score Max Prob & 64.00 & 60.33 \\
Z-Score Min Ent & 67.53 & 62.33 \\
Z-Score Energy & 71.67 & 62.80 \\
\bottomrule
\end{tabular}
\end{table}

\textbf{Analysis and Key Insights:}
The results demonstrate that achieving high-accuracy task routing is entirely possible without training any routing networks or backpropagating gradients. Standard \textbf{Min Entropy} achieves an impressive overall routing accuracy of \textbf{75.27\%}, leading to a final classification accuracy of \textbf{72.33\%} (recovering 88.5\% of the oracle performance). 

However, looking at the confusion matrices, standard metrics heavily favor simpler classification domains: both Max Probability and Min Entropy route MNIST and Fashion-MNIST with near-perfect accuracy ($>85\%$), but frequently misroute the more complex CIFAR-10 domain (routing only 232/500 CIFAR-10 images correctly). 
Our proposed \textbf{Calibrated Min Entropy} resolves this bias by normalizing test-time entropy against the expected calibration-set entropy. This yields the highest overall routing accuracy of \textbf{75.73\%} and distributes routing predictions much more evenly across all domains, correctly routing \textbf{387/500} CIFAR-10 images. 

Our newly introduced **Z-Score standardized metrics** systematically calibrate the scale of predictions by utilizing both the mean and standard deviation from the calibration set. Specifically, **Z-Score Energy** correctly routes **426/500** CIFAR-10 samples (85.2\% accuracy), behaving as an exceptional high-complexity domain detector. While standard Min Entropy remains the top-performer for downstream classification accuracy (72.33\%) due to its near-perfect accuracy on the simpler high-performance domains, the Calibrated and Z-Score standardized metrics provide a much more balanced and theoretically sound routing paradigm that prevents simpler domains from starving more complex ones.

These findings offer deep pragmatic value for practitioners. Rather than keeping multiple copies of massive backbones in memory or employing expensive, parameterized routers, developers can implement a simple, training-free entropy or energy-based router on top of TCAC. This unlocks a highly flexible and fully autonomous multi-task deployment pipeline with minimal computational overhead.

\section{Discussion, Limitations, and Future Work}
\label{sec:discussion_limitations}
While Task-Conditional Activation Calibration (TCAC) demonstrates outstanding empirical performance and rigorous theoretical guarantees, we must honestly evaluate its practical limitations and delineate avenues for future development.

\textbf{1. Dependency on BatchNorm Layers:} 
In its current formulation, TCAC leverages channel-wise normalization, which is naturally suited for models incorporating Batch Normalization (BatchNorm) layers (e.g., convolutional networks like ResNet or ConvNeXt). However, modern LLMs and vision transformers (ViTs) predominantly employ Layer Normalization (LayerNorm) or Root Mean Square Normalization (RMSNorm). Extending TCAC to those architectures would require calibrating token-wise or channel-wise activation statistics immediately before feedforward blocks, or developing a "LayerNorm-compatible" calibration scheme. This represents a highly promising direction for future research.

\textbf{2. Requirement for a Small Calibration Dataset:} 
Although TCAC is exceptionally data-efficient (requiring as few as 32 unlabeled samples per task, and remaining robust to extreme class-imbalance), it still requires access to a tiny subset of target-domain data. In zero-data scenarios where absolutely no target samples can be accessed due to extreme privacy or compliance restrictions, TCAC cannot compute its calibration parameters. Future work could investigate synthetic data generation or generative priors to compute activation statistics in a zero-shot, data-free manner.

\textbf{3. Routing Errors in Blind Deployments:}
In blind multi-task deployments where incoming samples lack task labels, our training-free routing framework (such as Calibrated Min Entropy or Z-Score Energy) achieves strong routing accuracy ($\sim 75.73\%$), but is not perfect. A misrouted sample uses the wrong calibration parameters and classification head, which limits downstream accuracy compared to the oracle setup. Future efforts could focus on combining training-free entropy/energy metrics with a lightweight, trainable parameter-free meta-classifier to further enhance routing robustness without sacrificing the zero-training paradigm.

\section{Conclusion}
\label{sec:conclusion}
We introduced Task-Conditional Activation Calibration (TCAC), a highly pragmatic, training-free approach to resolve representation shifts and variance collapse in multi-task model merging. While standard model merging methods suffer from significant parameter interference, and existing calibration methods like REPAIR fail by forcing a single averaged representation space, TCAC restores task-specific activation statistics using a tiny unlabeled calibration set. 

Through extensive evaluation, we demonstrated that TCAC is highly data-efficient (requiring as few as 32 samples per task), extremely robust to weight arithmetic hyperparameters, and offers a highly flexible depth-latency trade-off (with entry-layer calibration alone recovering over 80\% accuracy). Our work demonstrates that keeping model backbones shared while preserving task-conditional representation spaces is a powerful, elegant, and highly practical lever for post-merging alignment, unlocking high-performance multi-task models with zero training cost.

\bibliography{example_paper}
\bibliographystyle{icml2026}

%%%%%%%%%%%%%%%%%%%%%%%%%%%%%%%%%%%%%%%%%%%%%%%%%%%%%%%%%%%%%%%%%%%%%%%%%%%%%%%
%%%%%%%%%%%%%%%%%%%%%%%%%%%%%%%%%%%%%%%%%%%%%%%%%%%%%%%%%%%%%%%%%%%%%%%%%%%%%%%
% APPENDIX
%%%%%%%%%%%%%%%%%%%%%%%%%%%%%%%%%%%%%%%%%%%%%%%%%%%%%%%%%%%%%%%%%%%%%%%%%%%%%%%
%%%%%%%%%%%%%%%%%%%%%%%%%%%%%%%%%%%%%%%%%%%%%%%%%%%%%%%%%%%%%%%%%%%%%%%%%%%%%%%
\newpage
\appendix
\onecolumn
\section{Theoretical Justification of Multi-Task Variance Collapse and TCAC}
\label{sec:appendix_proof}

In this section, we provide a rigorous mathematical analysis of the multi-task variance collapse phenomenon observed in parameter-level model merging, and formally prove that our proposed Task-Conditional Activation Calibration (TCAC) restores the original representation statistics of the expert models.

\subsection{Multi-Task Variance Collapse Proof}
Let $\Phi_l(x)$ denote the intermediate representation (activations) at layer $l \in \{1, \dots, L\}$ of a deep neural network, and let $X_l^{(k)}$ represent the activations under the $k$-th task-specific expert model $\theta_k$. For simplicity, consider a sequence of $L$ linear or locally linear layers:
\begin{equation}
X_l^{(k)} = W_l^{(k)} X_{l-1}^{(k)}
\end{equation}
where $W_l^{(k)} \in \mathbb{R}^{d \times d}$ represents the weight matrix of expert $k$ at layer $l$. 

Suppose we merge the $K$ expert backbones at the weight level using simple weight averaging (WA):
\begin{equation}
W_l^{merged} = \frac{1}{K} \sum_{k=1}^K W_l^{(k)}
\end{equation}
To analyze the propagation of statistics, we make the following standard statistical assumptions:
\begin{assumption}
The expert weight matrices $W_l^{(k)}$ are independent and identically distributed across different tasks $k$, with mean $\mathbb{E}[W_l^{(k)}] = 0$ and element-wise variance $\text{Var}((W_l^{(k)})_{ij}) = \frac{\sigma_W^2}{d}$ for normalization.
\end{assumption}

Under the merged model, the uncalibrated activation propagates as:
\begin{equation}
X_l^{(merged)} = W_l^{merged} X_{l-1}^{(merged)} = \left(\frac{1}{K} \sum_{k=1}^K W_l^{(k)}\right) X_{l-1}^{(merged)}
\end{equation}
We now prove that the variance of the merged activations collapses exponentially relative to the expert models.
\begin{theorem}
Let the input variance be $\text{Var}(X_0^{(merged)}) = \text{Var}(X_0^{(k)}) = \sigma_0^2 \cdot I_d$. Under Assumption 3.2, the variance of the activation at layer $L$ under the uncalibrated merged model is:
\begin{equation}
\text{Var}(X_L^{(merged)}) = \left(\frac{1}{K}\right)^L \text{Var}(X_L^{(k)})
\end{equation}
\end{theorem}

\begin{proof}
We proceed by induction on the layer index $l$.
For the base case $l=1$, the activations are:
\begin{equation}
X_1^{(merged)} = W_1^{merged} X_0, \quad X_1^{(k)} = W_1^{(k)} X_0
\end{equation}
Since the input $X_0$ is identical and independent of the weights, we compute the covariance of $X_1^{(merged)}$:
\begin{align}
\text{Var}(X_1^{(merged)}) &= \mathbb{E}[ (W_1^{merged} X_0) (W_1^{merged} X_0)^T ] \\
&= \mathbb{E}\left[ \left(\frac{1}{K} \sum_{k=1}^K W_1^{(k)}\right) X_0 X_0^T \left(\frac{1}{K} \sum_{j=1}^K W_1^{(j)}\right)^T \right]
\end{align}
Since the weights $W_1^{(k)}$ and $W_1^{(j)}$ are independent for $k \neq j$ and have zero mean, the cross-terms vanish:
\begin{align}
\text{Var}(X_1^{(merged)}) &= \frac{1}{K^2} \sum_{k=1}^K \mathbb{E}[ W_1^{(k)} X_0 X_0^T (W_1^{(k)})^T ] \\
&= \frac{1}{K^2} \sum_{k=1}^K \text{Var}(X_1^{(k)})
\end{align}
Since the experts are identically distributed, $\text{Var}(X_1^{(k)}) = \Sigma_1$ is identical for all $k$. Thus:
\begin{equation}
\text{Var}(X_1^{(merged)}) = \frac{1}{K^2} \cdot K \Sigma_1 = \frac{1}{K} \text{Var}(X_1^{(k)})
\end{equation}
This establishes the induction base case. Now assume that at layer $l-1$, the relation holds:
\begin{equation}
\text{Var}(X_{l-1}^{(merged)}) = \left(\frac{1}{K}\right)^{l-1} \text{Var}(X_{l-1}^{(k)})
\end{equation}
For layer $l$, using the independence of weights $W_l^{(k)}$ across tasks and their independence from the activation $X_{l-1}^{(merged)}$, we have:
\begin{align}
\text{Var}(X_l^{(merged)}) &= \mathbb{E}[ W_l^{merged} X_{l-1}^{(merged)} (X_{l-1}^{(merged)})^T (W_l^{merged})^T ] \\
&= \text{Var}(W_l^{merged}) \cdot \mathbb{E}[ X_{l-1}^{(merged)} (X_{l-1}^{(merged)})^T ]
\end{align}
Using the element-wise variance of the merged weights $\text{Var}((W_l^{merged})_{ij}) = \frac{1}{K} \text{Var}((W_l^{(k)})_{ij})$, the output variance propagates as:
\begin{align}
\text{Var}(X_l^{(merged)}) &= \frac{1}{K} \text{Var}(W_l^{(k)}) \cdot \text{Var}(X_{l-1}^{(merged)}) \\
&= \frac{1}{K} \text{Var}(W_l^{(k)}) \cdot \left(\frac{1}{K}\right)^{l-1} \text{Var}(X_{l-1}^{(k)}) \\
&= \left(\frac{1}{K}\right)^l \left( \text{Var}(W_l^{(k)}) \text{Var}(X_{l-1}^{(k)}) \right) \\
&= \left(\frac{1}{K}\right)^l \text{Var}(X_l^{(k)})
\end{align}
By induction, the proof is complete for any depth $L$.
\end{proof}

This exponential decay explains why uncalibrated model merging collapses catastrophically: as depth $L$ increases, the representation variance goes to zero, rendering the features indistinguishable.

\subsection{Theoretical Guarantee of TCAC Restoration}
We now show that TCAC mathematically guarantees the restoration of the exact representation statistics.
During the calibration phase, we record:
\begin{equation}
\mu_{l, k}^{orig} = \mathbb{E}[X_l^{(k)}], \quad (\sigma_{l, k}^{orig})^2 = \text{Var}(X_l^{(k)})
\end{equation}
\begin{equation}
\mu_{l, k}^{merged} = \mathbb{E}[X_l^{(merged)}], \quad (\sigma_{l, k}^{merged})^2 = \text{Var}(X_l^{(merged)})
\end{equation}
During the inference phase, the uncalibrated activation $X_l \equiv X_l^{(merged)}$ is transformed via task-conditional rescaling:
\begin{equation}
\hat{X}_l = \frac{X_l - \mu_{l, k}^{merged}}{\sigma_{l, k}^{merged}} \cdot \sigma_{l, k}^{orig} + \mu_{l, k}^{orig}
\end{equation}

By the linearity of the expectation operator:
\begin{align}
\mathbb{E}[\hat{X}_l] &= \mathbb{E}\left[ \frac{X_l - \mu_{l, k}^{merged}}{\sigma_{l, k}^{merged}} \cdot \sigma_{l, k}^{orig} + \mu_{l, k}^{orig} \right] \\
&= \frac{\mathbb{E}[X_l] - \mu_{l, k}^{merged}}{\sigma_{l, k}^{merged}} \cdot \sigma_{l, k}^{orig} + \mu_{l, k}^{orig}
\end{align}
Since the test distribution is assumed to be locally consistent with the calibration distribution, we have $\mathbb{E}[X_l] = \mu_{l, k}^{merged}$. Substituting this in yields:
\begin{equation}
\mathbb{E}[\hat{X}_l] = 0 + \mu_{l, k}^{orig} = \mu_{l, k}^{orig}
\end{equation}

Next, we analyze the variance of the calibrated activations:
\begin{align}
\text{Var}(\hat{X}_l) &= \text{Var}\left( \frac{X_l - \mu_{l, k}^{merged}}{\sigma_{l, k}^{merged}} \cdot \sigma_{l, k}^{orig} + \mu_{l, k}^{orig} \right) \\
&= \text{Var}\left( \frac{X_l}{\sigma_{l, k}^{merged}} \cdot \sigma_{l, k}^{orig} \right) \\
&= \frac{\text{Var}(X_l)}{(\sigma_{l, k}^{merged})^2} \cdot (\sigma_{l, k}^{orig})^2
\end{align}
Since the test activation variance matches the calibration variance $\text{Var}(X_l) = (\sigma_{l, k}^{merged})^2$, this simplifies to:
\begin{equation}
\text{Var}(\hat{X}_l) = 1 \cdot (\sigma_{l, k}^{orig})^2 = (\sigma_{l, k}^{orig})^2
\end{equation}

Thus, TCAC exactly restores both the first and second moments of the original expert activations at every calibrated layer $l$, preventing variance collapse and representation shift, thereby theoretically justifying its outstanding empirical performance.

\subsection{Empirical Robustness under Input Corruptions}
\label{sec:appendix_noise}

To evaluate our framework under realistic deployment challenges, we assess the robustness of both the uncalibrated baseline (Task Arithmetic with $\lambda=0.4$) and our proposed TCAC (Task-Specific Affine) under input-level image corruptions. Specifically, we evaluate two common real-world corruptions on the test sets of MNIST, Fashion-MNIST, and CIFAR-10: (1) \textbf{Gaussian Noise} with standard deviation $\sigma_{noise} \in \{0.0, 0.05, 0.10, 0.15, 0.20\}$, and (2) \textbf{Gaussian Blur} with blur kernel standard deviation $\sigma_{blur} \in \{0.0, 0.5, 1.0, 1.5, 2.0\}$. Crucially, consistent with a pragmatic deployment paradigm, our TCAC parameters are calibrated once on \textit{clean} calibration samples ($N=128$), while evaluation is performed on the corrupted test inputs.

The results are summarized in Table \ref{tab:noise_robustness_data} and visualized in Figure \ref{fig:noise_blur_plots}.

\begin{table}[h]
\centering
\caption{Model merging accuracies (\%) under input Gaussian Noise and Gaussian Blur corruptions across different task datasets (MNIST, Fashion-MNIST, CIFAR-10). The uncalibrated baseline performs at random guessing (approx.\ 10.0\%) across all corruption levels due to severe parameter interference, whereas our TCAC method degrades gracefully even under extreme noise and blur.}
\label{tab:noise_robustness_data}
\vskip 0.1in
\begin{tabular}{lcccccccc}
\toprule
 & \multicolumn{4}{c}{\textbf{Gaussian Noise Std ($\sigma_{noise}$)}} & \multicolumn{4}{c}{\textbf{Gaussian Blur Sigma ($\sigma_{blur}$)}} \\
\cmidrule(lr){2-5} \cmidrule(lr){6-9}
\textbf{Task \& Method} & \textbf{0.00} & \textbf{0.05} & \textbf{0.10} & \textbf{0.20} & \textbf{0.50} & \textbf{1.00} & \textbf{1.50} & \textbf{2.00} \\
\midrule
\textit{MNIST} \\
~~Baseline & 9.80 & 9.80 & 9.80 & 9.80 & 9.80 & 9.80 & 9.80 & 9.80 \\
~~TCAC & \textbf{96.87} & \textbf{96.48} & \textbf{95.05} & \textbf{81.33} & \textbf{97.03} & \textbf{95.60} & \textbf{86.84} & \textbf{61.05} \\
\midrule
\textit{Fashion-MNIST} \\
~~Baseline & 10.00 & 10.00 & 10.00 & 10.00 & 10.00 & 10.00 & 10.00 & 10.00 \\
~~TCAC & \textbf{86.41} & \textbf{85.33} & \textbf{80.76} & \textbf{60.79} & \textbf{86.01} & \textbf{76.11} & \textbf{59.67} & \textbf{45.54} \\
\midrule
\textit{CIFAR-10} \\
~~Baseline & 10.00 & 10.00 & 10.00 & 10.00 & 10.00 & 10.00 & 10.00 & 10.00 \\
~~TCAC & \textbf{63.46} & \textbf{62.69} & \textbf{58.06} & \textbf{45.66} & \textbf{59.30} & \textbf{40.07} & \textbf{23.30} & \textbf{14.37} \\
\midrule
\textit{Average} \\
~~Baseline & 9.93 & 9.93 & 9.93 & 9.93 & 9.93 & 9.93 & 9.93 & 9.93 \\
~~TCAC & \textbf{82.25} & \textbf{81.50} & \textbf{77.96} & \textbf{62.59} & \textbf{80.78} & \textbf{70.59} & \textbf{56.60} & \textbf{40.32} \\
\bottomrule
\end{tabular}
\end{table}

\begin{figure}[h]
    \centering
    \begin{subfigure}{0.49\textwidth}
        \centering
        \includegraphics[width=\textwidth]{plots/noise_robustness.png}
        \caption{Gaussian Noise Robustness}
        \label{fig:noise_robustness}
    \end{subfigure}
    \hfill
    \begin{subfigure}{0.49\textwidth}
        \centering
        \includegraphics[width=\textwidth]{plots/blur_robustness.png}
        \caption{Gaussian Blur Robustness}
        \label{fig:blur_robustness}
    \end{subfigure}
    \caption{Model performance sweeps across input-level Gaussian Noise and Gaussian Blur. While the uncalibrated baseline is entirely non-functional due to internal representation collapse, TCAC exhibits outstanding resilience, degrading smoothly even under extreme noise ($\sigma_{noise}=0.20$) or heavy blurring ($\sigma_{blur}=2.00$).}
    \label{fig:noise_blur_plots}
\end{figure}

\textbf{Analysis of Noise Robustness:}
The empirical results reveal that our TCAC framework is exceptionally resilient to input-level perturbations. Under Gaussian Noise, the uncalibrated baseline is completely non-functional across all noise levels (9.93\% average accuracy). In contrast, TCAC maintains a high average accuracy of \textbf{81.50\%} under mild noise ($\sigma_{noise}=0.05$) and only drops to \textbf{77.96\%} under moderate noise ($\sigma_{noise}=0.10$). Even under extreme noise ($\sigma_{noise}=0.20$), where the raw signal is severely corrupted, TCAC remains functional with an average accuracy of \textbf{62.59\%} (including 81.33\% on MNIST). This indicates that the calibrated feature statistics provide a stabilizing boundary that prevents the network from suffering from catastrophic out-of-distribution hallucinations under random weight perturbations.

\textbf{Analysis of Blur Robustness:}
Under Gaussian Blur, we observe a similar trend of graceful and smooth degradation. While the baseline remains at random guessing, TCAC retains high performance under mild blurring ($\sigma_{blur}=0.50$, average accuracy \textbf{80.78\%}) and moderate blurring ($\sigma_{blur}=1.00$, average accuracy \textbf{70.59\%}). As expected, natural images with rich high-frequency textures (CIFAR-10) are more sensitive to blurring, dropping to 14.37\% accuracy under heavy blurring ($\sigma_{blur}=2.00$), whereas simpler high-contrast grayscale objects (MNIST and Fashion-MNIST) remain highly recognizable, with TCAC achieving 61.05\% and 45.54\% respectively. 

These results demonstrate the outstanding practical utility and robustness of TCAC. In real-world edge applications, input sensors are frequently subjected to noise, defocus blur, or lens degradation. The fact that a training-free, zero-overhead representation alignment method like TCAC not only heals internal parameter conflicts but also preserves a high degree of out-of-distribution robustness to input corruptions makes it a highly compelling, reliable, and pragmatic choice for deploying merged models in the wild.

\subsection{Data-Free and Out-of-Domain Calibration Transferability}
\label{sec:appendix_data_free}

In Section~\ref{sec:discussion_limitations}, we identified the requirement of a small calibration set from the target domain as a limitation of our proposed TCAC framework. To address this bottleneck from a pragmatic perspective, we systematically investigate the transferability and resilience of TCAC under two challenging, data-scarce settings:
\begin{enumerate}
    \item \textbf{Data-Free Calibration}: We calibrate the merged model using completely synthetic, random Gaussian noise $X \sim \mathcal{N}(0, I)$ ($N=128$), introducing zero reliance on real data.
    \item \textbf{Cross-Domain Out-of-Domain Calibration}: We calibrate the merged model using completely out-of-domain real images. Specifically, we calibrate MNIST and Fashion-MNIST models using CIFAR-10 images, and we calibrate the CIFAR-10 model using MNIST images ($N=128$).
\end{enumerate}
Evaluating under these settings allows us to assess whether intermediate activation statistics represent domain-specific semantics or if they primarily reflect structural weight-space properties that can be extracted using out-of-domain signals.

The downstream classification accuracies on the test sets of MNIST, Fashion-MNIST, and CIFAR-10 are presented in Table~\ref{tab:data_free_results} and visualized in Figure~\ref{fig:data_free_plot}.

\begin{table}[h]
\centering
\caption{Model merging accuracies (\%) under Data-Free (Gaussian Noise) and Cross-Domain (Out-of-Domain) calibration settings compared to clean in-domain calibration and the uncalibrated baseline (Task Arithmetic, $\lambda=0.4$, $N=128$). Utilizing real out-of-domain images for calibration achieves spectacular performance recovery, proving that TCAC statistics are highly transferable.}
\label{tab:data_free_results}
\vskip 0.1in
\begin{tabular}{lcccc}
\toprule
\textbf{Calibration Setting} & \textbf{MNIST} & \textbf{Fashion-MNIST} & \textbf{CIFAR-10} & \textbf{Average} \\
\midrule
Uncalibrated Baseline & 9.80 & 10.00 & 10.00 & 9.93 \\
\midrule
Data-Free (Gaussian Noise) & 11.35 & 10.00 & 13.00 & 11.45 \\
Cross-Domain (Out-of-Domain) & \textbf{72.78} & \textbf{76.78} & \textbf{23.57} & \textbf{57.71} \\
Clean In-Domain (Oracle) & 96.92 & 86.47 & 63.65 & 82.35 \\
\bottomrule
\end{tabular}
\end{table}

\begin{figure}[h]
    \centering
    \includegraphics[width=0.6\textwidth]{plots/data_free_robustness.png}
    \caption{TCAC performance comparison under Data-Free (Gaussian Noise) and Cross-Domain (Out-of-Domain) calibration schemes. While purely random noise fails to excite the network layers constructively, utilizing real out-of-domain natural images (e.g., CIFAR-10) for grayscale calibration recovers up to 76.78\% accuracy, demonstrating outstanding calibration transferability.}
    \label{fig:data_free_plot}
\end{figure}

\textbf{Failure of Synthetic Gaussian Noise}:
The empirical results show that calibrating TCAC using purely synthetic Gaussian noise fails to recover performance (average accuracy of 11.45\% vs.\ 9.93\% baseline). This occurs because white noise lacks the spatial coherence and local features characteristic of real images, causing it to propagate through the deep network differently. Consequently, the recorded statistics are misaligned with the natural image manifolds, making them ineffective for calibrating test-time activations.

\textbf{Spectacular Transferability of Cross-Domain Calibration}:
In stark contrast, calibrating TCAC using completely out-of-domain real images achieves an outstanding and unexpected recovery. Calibrating the MNIST model using CIFAR-10 natural images recovers the classification accuracy from 9.80\% baseline to \textbf{72.78\%}. Similarly, calibrating the Fashion-MNIST model using CIFAR-10 images boosts the performance from 10.00\% baseline to \textbf{76.78\%}. Even the complex CIFAR-10 model, when calibrated with simple grayscale MNIST digits, improves from 10.00\% baseline to \textbf{23.57\%}. On average across tasks, Cross-Domain Calibration recovers \textbf{57.71\%} accuracy, closing a major portion of the gap between the uncalibrated baseline and the clean oracle.

\textbf{Pragmatic Implications for Real-World Deployment}:
This is a highly significant and encouraging result for practitioners. In real-world edge deployments, obtaining unlabeled data from the target domain can often be extremely difficult due to operational, privacy, or safety constraints (e.g., medical imaging or confidential industrial inspections). Our findings empirically demonstrate that developers do not need target-domain data to perform successful activation calibration. By using \textit{any readily available real-world dataset} (such as ImageNet or CIFAR-10), they can capture the network's structural layer-wise variance collapse behavior and calibrate the shared backbone, restoring over 70\% of the individual expert performance with zero target data. This proves that TCAC is highly transferable, robust, and exceptionally well-suited for blind, resource-scarce environments.

\end{document}
